% Part: incompleteness
% Chapter: representability-in-q
% Section: undecidability

\documentclass[../../../include/open-logic-section]{subfiles}

\begin{document}

\olfileid[fr]{inc}{req}{und}
\olsection{Indécidabilité}

Une théorie~$\Th{T}$ est dite \emph{indécidable} s'il n'existe
aucune procédure de calcul qui, en un nombre fini d'étapes,
réponde toujours correctement à la question « $\Th{T}$
démontre-t-elle~$!A$ ? » pour tout !!{sentence}~$!A$ de son
langage. Ainsi, $\Th{Q}$ serait décidable s'il existait une
procédure qui, à partir d'un !!{sentence}~$!A$ du langage de
l'arithmétique, décide si $\Th{Q}\Proves !A$. Plus
précisément : la relation $\Prov[\Th{Q}](y)$, vraie si et
seulement si $y$ est le numéro de Gödel d'un !!{sentence}
démontrable dans~$\Th{Q}$, est-elle récursive ? Non.

\begin{thm}
$\Th{Q}$ est indécidable : la relation
\[
\Prov[\Th{Q}](y) \defiff \fn{Sent}(y) \land
\lexists[x][\Prf[\Th{Q}](x, y)]
\]
n'est pas récursive.
\end{thm}

\begin{proof}
Supposons qu'elle le soit. Nous pourrions alors résoudre le
problème de l'arrêt. Étant donnés $e,n$, nous savons que
$\cfind{e}(n)\fdefined$ si et seulement s'il existe un~$s$
tel que $T(e,n,s)$, où $T$ est le prédicat de Kleene du
\olref[cmp][rec][nft]{thm:kleene-nf}. Puisque $T$ est
récursif primitif, une !!{formula}~$!B_T$ le représente
dans~$\Th{Q}$. Ainsi, si $T(e,n,s)$ est vraie, alors
$\Th{Q}\Proves !B_T(\num e,\num n,\num s)$; si elle est
fausse, $\Th{Q}\Proves\lnot !B_T(\num e,\num n,\num s)$.
La cohérence de $\Th{Q}$ donne également
$\Th{Q}\Proves !B_T(\num e,\num n,\num s)$ si et seulement
si $T(e,n,s)$. Une instance positive permet de conclure
$\Th{Q}\Proves\lexists[y][!B_T(\num e,\num n,y)]$.

Si aucun~$s$ ne satisfait $T(e,n,s)$, alors
$\Th{Q}\Proves\lnot !B_T(\num e,\num n,\num s)$
pour tout~$s$. Or $\Th{Q}$ est $\omega$-cohérente : si,
pour chaque $n\in\Nat$,
$\Th{Q}\Proves\lnot !A(\num n)$, alors
$\Th{Q}\Proves/\lexists[y][!A(y)]$.
Cela découle du fait que les axiomes de $\Th{Q}$ sont vrais
dans le modèle standard~$\Struct{N}$. Donc
$\Th{Q}\Proves/\lexists[y][!B_T(\num e,\num n,y)]$.
En somme, $\Th{Q}$ démontre cette existentiale si et
seulement s'il existe un~$s$ tel que $T(e,n,s)$, autrement
dit si et seulement si $\cfind{e}(n)\fdefined$.

À partir de $e,n$, on peut calculer
$\Gn{\lexists[y][!B_T(\num e,\num n,y)]}$; notons
$g(e,n)$ la fonction récursive primitive qui le fait.
Alors
\[
h(e, n) =
\begin{cases}
1 & \text{si $\Prov[\Th{Q}](g(e, n))$}\\
0 & \text{sinon}.
\end{cases}
\]
Si $\Prov[\Th{Q}]$ était récursive, $h$ le serait
aussi. Mais \olref[cmp][rec][hlt]{thm:halting-problem}
établit que $h$ ne l'est pas. La relation
$\Prov[\Th{Q}]$ ne peut donc pas être récursive.
\end{proof}

\begin{cor}
La logique du premier ordre est indécidable.
\end{cor}

\begin{proof}
Si la logique du premier ordre était décidable, la
démontrabilité dans~$\Th{Q}$ le serait aussi, puisque
$\Th{Q}\Proves !A$ si et seulement si
$\Proves !T\lif !A$, où $!T$ est la conjonction des axiomes
de~$\Th{Q}$.
\end{proof}

\end{document}
